\question{}

\begin{ابجد}
\item طبق رابطه \lr{Policy Evaluation} ارزش حالت‌ها را به‌روز می‌کنیم:
\begin{align*}
    V_{k+1}(s) = \sum_{s'}T(s, \pi(s), s')[R(s, \pi(s), s') + \gamma V_k(s')]
\end{align*}
اجرای این بخش به این صورت می‌باشد:
\begin{align*}
     & V_2(A) = \frac{1}{2}(4 + 0.5(2)) + \frac{1}{2}(2 + 0.5(2)) = 4                                          \\
     & V_2(B) = \frac{1}{2}(0 + 0.5(2)) + \frac{1}{2}(2 + 0.5(2)) = 2                                          \\
     & V_2(C) = \frac{1}{2}(1 + 0.5(2)) + \frac{1}{2}(\frac{1}{4}(8 + 0.5(2)) + \frac{3}{4}(8 + 0.5(2))) = 5.5
\end{align*}

\item حال باید طبق رابطه‌ی زیر، \lr{Policy Improvement} را انجام دهیم:
\begin{align*}
    \pi_{i + 1}(s) = \operatorname*{argmax}_a \sum_{s'}T(s, a, s')[R(s, a, s') + \gamma V^{\pi_i}(s')]
\end{align*}
بنابراین در این بخش داریم:
\begin{align*}
     & s=A &  & a=x: 1(4 + 0.5(2)) = 5                                                                        \\
     &     &  & a=y: 1(2 + 0.5(5.5)) = 4.75                                        & \Rightarrow \pi_2(A) = x \\
    \\
     & s=B &  & a=x: 1(2 + 0.5(5.5)) = 4.75                                                                   \\
     &     &  & a=y: 1(0 + 0.5(4)) = 2                                             & \Rightarrow \pi_2(B) = x \\
    \\
     & s=C &  & a=x: \frac{1}{4}(8 + 0.5(4)) + \frac{3}{4}(8 + 0.5(5.5)) = 10.5625                            \\
     &     &  & a=y: 1(1 + 0.5(2)) = 2                                             & \Rightarrow \pi_2(C) = x \\
\end{align*}

\item ابتدا نشان می‌دهیم اگر $\pi'$ به صورت حریصانه از $\pi$ به دست آمده باشد، می‌توان نتیجه گرفت به ازای تمام حالت‌ها داریم $V^{\pi'}(s) \geq V^{\pi}(s)$.
\begin{proof}
    طبق تعریف داریم:
    \begin{align*}
                          & Q^{\pi}(s, \pi'(s)) = \max_a Q^{\pi}(s, a)  \\
        \Rightarrow \quad & Q^{\pi}(s, \pi'(s)) \geq Q^{\pi}(s, \pi(s)) \\
        \Rightarrow \quad & Q^{\pi}(s, \pi'(s)) \geq V^{\pi}(s)
    \end{align*}
    همچنین می‌دانیم $Q^{\pi}(s, \pi'(s)) = \mathbb{E}_{\pi'} [R_{t+1} + \gamma V^\pi(S_{t+1})]$، بنابراین:
    \begin{align*}
        V^{\pi}(s) \leq \mathbb{E}_{\pi'} [R_{t+1} + \gamma V^{\pi}(s_{t+1})] \quad \Rightarrow \quad V^{\pi}(s) \leq \mathbb{E}_{\pi'} [R_{t+1} + \gamma Q^{\pi}(s_{t+1}, \pi'(s_{t+1}))]
    \end{align*}
    و می‌توانیم این فرایند را به طور زنجیره‌ای ادامه دهیم:
    \begin{align*}
         & V^{\pi}(s) \leq \mathbb{E}_{\pi'} [R_{t+1} + \gamma R_{t+2} + \gamma^2 V^{\pi}(s_{t+2})]                               \\
         & \vdots                                                                                                                 \\
         & V^{\pi}(s) \leq \mathbb{E}_{\pi'} [R_{t+1} + \gamma R_{t+2} + \dots + \gamma^{n-1} R_{t+n} +\gamma^n V^{\pi}(s_{t+n})]
    \end{align*}
    از آنجا که $0 \leq \gamma < 1$، با میل کردن $n \to \infty$، عبارت $\gamma^n V^{\pi}(s_{t+n})$ به صفر میل کرده و تنها عبارت‌های تابع پاداش باقی می‌مانند. به عبارتی:
    \begin{align*}
        V^\pi(s) \le \mathbb{E}_{\pi'} [ \sum_{k=0}^{\infty} \gamma^k R_{t+k+1}] = V^{\pi'}(s)
    \end{align*}
    پس نشان دادیم به ازای هر حالت دلخواه، داریم $V^{\pi'}(s) \geq V^{\pi}(s)$.
\end{proof}

حالا باید نشان دهیم که اگر حالت تساوی رخ دهد، سیاست $\pi'$ بهینه است.
\begin{proof}
    اگر در روابط قسمت قبل تساوی قرار دهیم، داریم:
    \begin{align*}
        V^{\pi'}(s) = \max_a Q^{\pi}(s, a) = \max_a \sum_{s'}T(s, a, s')[R(s,a,s') + \gamma V^{\pi}(s')]
    \end{align*}
    همچنین فرض کردیم $V^{\pi}(s) = V^{\pi'}(s)$. بنابراین:
    \begin{align*}
        V^{\pi'}(s) = \max_a \sum_{s'}T(s, a, s')[R(s,a,s') + \gamma V^{\pi'}(s')]
    \end{align*}
    می‌بینیم که $V^{\pi'}(s)$ در رابطه‌ی بلمن صدق می‌کند. بنابراین $\pi'$ سیاست بهینه است.
\end{proof}

\end{ابجد}